(a) On $x_i\ne0$, the coordinate ring is $B=(S(Y)_{x_i})_0$. If an element of $\operatorname{Frac}(B)$ is integral over $B$, it lies in the normal ring $S(Y)_{x_i}$. Multiplying it by a nonzero degree-zero denominator shows, by the grading, that it has degree zero. Thus it belongs to $B$, so the affine charts and hence $Y$ are normal.

(b) Here $S(Y)=k[t^4,t^3u,tu^3,u^4]$. The element $t^2u^2=(t^3u)^2/t^4$ belongs to its fraction field and is integral, since its square is $t^4u^4$. It does not belong to $S(Y)$: the degree-four part is spanned by the four displayed monomials. Thus $Y$ is not projectively normal. It is normal by (c).

(c) The inverse to $(t:u)\mapsto(t^4:t^3u:tu^3:u^4)$ is $(x:y:z:w)\mapsto(x:y)$ on $x\ne0$, and $(z:w)$ on $w\ne0$. These charts cover $Y$, and the maps agree on their overlap. Thus $Y\cong\mathbb P^1$, whose homogeneous coordinate ring $k[t,u]$ is integrally closed.
